%------------------------------------------------------------------------------------
%  TABLEAU 1 : SYNTHÈSE DE MESURES EXPÉRIMENTALES (booktabs + tabularx + siunitx)
%------------------------------------------------------------------------------------
\begin{table}[htbp]
  \centering
  \small
  \caption{Synthèse des mesures thermomécaniques et validation expérimentale}
  \label{tab:synthese_mesures}
  \begin{tabularx}{\textwidth}{l c c c S[table-format=3.2] X}
    \toprule
    \textbf{Échantillon} & \textbf{Température} & \textbf{Pression} & \textbf{Vitesse} & {\textbf{Rendement (\%)}} & \textbf{Observations qualitatives} \\
    \midrule
    ECH-01 & \qty{22.4}{\degreeCelsius} & \qty{101.3}{\kilo\pascal} & \qty{1.45}{\meter\per\second} & 88.50 & Écoulement laminaire stable sans turbulence. \\
    ECH-02 & \qty{45.0}{\degreeCelsius} & \qty{150.2}{\kilo\pascal} & \qty{2.80}{\meter\per\second} & 92.15 & Amorçage de vortex en sortie de tuyère. \\
    ECH-03 & \qty{75.8}{\degreeCelsius} & \qty{210.0}{\kilo\pascal} & \qty{4.12}{\meter\per\second} & 95.40 & Régime turbulent pleinement développé. \\
    \addlinespace
    ECH-04 & \qty{90.2}{\degreeCelsius} & \qty{280.5}{\kilo\pascal} & \qty{5.60}{\meter\per\second} & 94.80 & Perte de charge par friction pariétale. \\
    ECH-05 & \qty{110.0}{\degreeCelsius} & \qty{350.0}{\kilo\pascal} & \qty{7.25}{\meter\per\second} & 91.20 & cavitation détectée au col d'étranglement. \\
    \bottomrule
  \end{tabularx}
\end{table}

